\subsection{Máquina de Estados Finita}

Uma máquina de estados é um modelo matemático utilizado para descrever sistemas que evoluem ao longo do tempo de acordo com entradas e estados internos. Em termos de circuitos digitais, ela representa formalmente sistemas sequenciais, nos quais a saída depende não apenas das entradas atuais, mas também do estado armazenado previamente em memória.

Uma Máquina de Estados Finita (FSM - Finite State Machine) é um modelo no qual o número de estados possíveis é limitado. Cada estado representa uma configuração discreta do sistema, e as transições entre estados são definidas por regras lógicas baseadas nas entradas e no estado atual.

Formalmente, uma FSM pode ser descrita por uma quíntupla:

\[
FSM = (Q, \Sigma, \delta, q_0, F)
\]

onde:
\begin{itemize}
    \item $Q$ é o conjunto finito de estados
    \item $\Sigma$ é o conjunto de entradas
    \item $\delta$ é a função de transição de estados
    \item $q_0$ é o estado inicial
    \item $F$ é o conjunto de estados finais (opcional em sistemas cíclicos)
\end{itemize}

\subsubsection{Cronômetro como Máquina de Estados}

No contexto deste projeto, o cronômetro pode ser modelado como uma máquina de estados finita composta por três subsistemas sequenciais encadeados, representando segundos, minutos e horas.

Cada estado global do sistema pode ser definido como uma tupla:

\[
\mathcal{S} = (H, M, S, D)
\]

onde:

\begin{align*}
H &\in [0, 23] \quad \text{representa as horas} \\
M &\in [0, 59] \quad \text{representa os minutos} \\
S &\in [0, 59] \quad \text{representa os segundos} \\
D &\in \{0,1\} \quad \text{representa a direção (0 = crescente, 1 = decrescente)}
\end{align*}

Como será detalhado posteriormente, os elementos $M$ e $S$ são implementados internamente como pares de dígitos decimais, divididos em dezena e unidade, tais que:

\[
M = (M_{dezena}, M_{unidade}), \quad S = (S_{dezena}, S_{unidade})
\]

com:

\[
M_{dezena}, S_{dezena} \in [0,5], \quad
M_{unidade}, S_{unidade} \in [0,9]
\]

O conjunto de estados válidos do sistema é, portanto, finito e dado por:

\[
|\mathcal{S}| = 24 \cdot 60 \cdot 60 \cdot 2
\]

Seja  $S_t$  o estado global do sistema em determinado momento $t$, a função de transição de estados pode ser descrita genericamente por:

\[
\delta(\mathcal{S}_t) = \mathcal{S}_{t+1}
\]

ou, de forma expandida:

\[
\delta(H_t, M_t, S_t, D_t) = (H_{t+1}, M_{t+1}, S_{t+1}, D_{t+1})
\]

Retornaremos a esta definição mais à frente no documento.